\section{Conclusiones y trabajo futuro}

Se completó un simulador del modelo extendido de Schelling que permite repetir una ejecución con los mismos resultados. La versión secuencial sirve como referencia y la versión paralela reparte franjas de la grilla con MPI y usa OpenMP para calcular satisfacción, precios y destinos. Ambas versiones produjeron el mismo estado con distintas cantidades de procesos e hilos. También se comprobó que no se pierdan hogares ni viviendas vacías y que una ejecución guardada pueda continuar correctamente.

Las mediciones mostraron que OpenMP acelera la búsqueda cuando existe suficiente trabajo. En la grilla principal, la versión secuencial demoró 3~h 49~min y la mejor configuración, con un proceso y cuatro hilos, demoró 2~h 6~min: una mejora de 1,817. Dos hilos lograron 1,724 con mayor eficiencia. Agregar una segunda máquina no mejoró el mejor tiempo porque cada proceso conserva datos globales y las configuraciones de dos nodos intercambiaron aproximadamente 83,2~GB.

La grilla principal usa permanencia cero, financiación de 60\% y 20\% de viviendas vacías. La satisfacción pasó de 54,6\% a 78,4\% y se aceptaron casi 2,5 millones de mudanzas. La proporción media de vecinos de la misma clase aumentó de 43,7\% a 59,3\%, y la grilla final mostró agrupamientos más claros. Estos resultados coinciden con el comportamiento esperado del modelo de Schelling, dentro de las limitaciones de una grilla sintética.

\subsection{Trabajo futuro}

La versión MPI podría repartir las viviendas vacías entre los procesos en lugar de mantener una copia completa en cada uno. Esto reduciría la memoria y los datos intercambiados, aunque obligaría a coordinar la búsqueda y los conflictos entre procesos. También sería útil repetir la grilla principal con varias semillas y varias ejecuciones por configuración. Finalmente, usar un mapa y datos censales por zona permitiría representar Montevideo con mayor detalle.
